\chapter{Literature Review}

\section{Classical Chess Engines}
% Minimax algorithm and alpha-beta pruning
% Evaluation functions and their limitations
% Representative engines (Stockfish architecture)

\section{Probabilistic Search Methods}
% Monte Carlo Tree Search (MCTS) principles
% MCTS in game playing
% Advantages over minimax

\section{The Deep Learning Revolution}
% AlphaZero breakthrough [Silver et al., 2017]
% Architecture: policy and value networks
% Self-play reinforcement learning
% Leela Chess Zero as open-source implementation

\section{Supervised Learning Approaches}
% DeepChess [David et al., 2017]
% Giraffe [Lai, 2015]
% Learning from human games vs. self-play

\section{Recent Innovations}
% Graph representations [Flet-Berliac & D'Anvers, 2024]
% Explainable AI in chess [Svendsen & Tørresen, 2022]
% Hybrid approaches

\section{Game AI Frameworks}
% OpenSpiel architecture and capabilities
% Python-chess library
% Related frameworks

\section{Web-Based Chess Platforms}
\subsection{Existing Interactive Chess Platforms}
% Lichess, Chess.com bot arenas
% Browser-based game interfaces
\subsection{Bot-vs-Bot Comparison Platforms}
% Existing tools for pitting AI agents against each other
% Spectator and analysis features
\subsection{Human-vs-Bot Interfaces}
% Existing platforms allowing users to play against AI
% UX considerations for interactive play
\subsection{Web Technologies for Real-Time Games}
% WebSockets, server-side game orchestration
% Frontend frameworks for board rendering

\section{Critical Analysis}
% Strengths and weaknesses of each approach
% Gaps in current research
% How your project builds on/differs from existing work

\section{Summary}
% Key insights informing your methodology
